% STATUS: written 2026-11-06 (English gate clean; sources read in full as recorded below).
% Sources read in full: Weibel 1994 (sections 9.1, 9.2, 9.4, 9.5),
% Atesli--Esen--Sutlu 2021 (pp. 1--3, 25--26), Xu--Liu 2021 (pp. 1--3),
% Das 2025 (pp. 1--3, 6--7, 19--20).
% Sources: Weibel1994, AtesliEsenSutlu2021, XuLiu2021ThreeLieDer, Das2025Nijenhuis.
\section{Computations for classes of algebras}\label{sec:cohomology:computations}

Each computation below is a derivation or a cohomology class of one, on a named class of algebras.

The theory is computable on the standard classes, and Weibel collects the
answers~\cite[\S 9]{Weibel1994}.

The tensor algebra is the first class, and its Hochschild homology vanishes
above dimension one (Proposition~\ref{prop:cohomology:tensor})
\cite[Proposition~9.1.6, pp.~303--304]{Weibel1994}.\begin{proposition}[{\cite[Proposition 9.1.6, pp.~303--304]{Weibel1994}}]\label{prop:cohomology:tensor}
Let $T = T(V)$ be the tensor algebra of a $k$-module $V$, and let $M$ be a
$T$-$T$ bimodule. Then $H_i(T, M) = 0$ for $i \neq 0, 1$, and there is an exact
sequence
\[
0 \longrightarrow H_1(T, M) \longrightarrow M \otimes V \xrightarrow{\;b\;} M
\longrightarrow H_0(T, M) \longrightarrow 0,
\]
where $b$ is the map $b(m \otimes v) = mv - vm$.
\end{proposition}

Polynomial rings sit inside this class~\cite[Exercise~9.1.3, p.~304]{Weibel1994}.
For $R = k[x_1, \dots, x_n]$ one has $H_i(R, M) = 0$ for $i > 0$ over a
field~\cite[Exercise~9.1.3, p.~304]{Weibel1994}. Weibel marks this as a special
case of Theorem~\ref{thm:cohomology:hkr} below.

For a commutative algebra the first group is built from the differentials of
Section~\ref{sec:commutative:kahler}.

\begin{proposition}[{\cite[Proposition 9.2.2, pp.~306--307]{Weibel1994}}]\label{prop:cohomology:comm-h1}
Let $R$ be a commutative $k$-algebra, and let $M$ be a right $R$-module made into
an $R$-$R$ bimodule by $rm = mr$. Then $H_0(R, M) \cong M$ and
$H_1(R, M) \cong M \otimes_R \Omega_{R/k}$.
\end{proposition}

The Hochschild--Kostant--Rosenberg theorem computes the whole graded theory for
a smooth commutative algebra~\cite[Theorem~9.4.7, pp.~322--323]{Weibel1994}. Its
Hochschild homology is then the exterior algebra on the Kähler differentials.

\begin{theorem}[{Hochschild--Kostant--Rosenberg, \cite[Theorem 9.4.7, pp.~322--323]{Weibel1994}}]\label{thm:cohomology:hkr}
Let $R$ be a commutative algebra, essentially of finite type over a field $k$. If
$R$ is smooth over $k$, then $\psi$ is an isomorphism of graded $R$-algebras
$H_*(R, R) \cong \Lambda^* \Omega_{R/k}$.
\end{theorem}

The degree one case of Theorem~\ref{thm:cohomology:hkr} gives
$H_1(R, R) \cong \Omega_{R/k}$~\cite[Theorem~9.4.7, p.~322]{Weibel1994}. That
is Proposition~\ref{prop:cohomology:comm-h1} read on $M = R$, and it feeds the
de Rham comparison of
Section~\ref{sec:cohomology:algebraic-topology}
\cite[Theorem~9.8.13, p.~351]{Weibel1994}. The derivation side of the same
degree is the universal derivation of
Section~\ref{sec:commutative:kahler}. % claim-ok: cross-reference sentence, the module is built there

Morita equivalence leaves the theory
unchanged~\cite[Theorem~9.5.6, p.~325]{Weibel1994}. A theorem of R.~K.~Dennis
says that this unchanged value holds across the
equivalence~\cite[Theorem~9.5.6, p.~325]{Weibel1994}, and it comes next (Theorem~\ref{thm:cohomology:morita}).

\begin{theorem}[{R.~K.~Dennis, \cite[Theorem 9.5.6, p.~325]{Weibel1994}}]\label{thm:cohomology:morita}
Hochschild homology is Morita invariant: if $R$ and $S$ are Morita equivalent
rings and $M$ is an $R$-$R$ bimodule, then
$H_*(R, M) \cong H_*(S, Q \otimes_R M \otimes_R P)$.
\end{theorem}

So the Hochschild theory of a matrix ring agrees with that of the ring
itself~\cite[Theorem~9.5.6, p.~325]{Weibel1994}.

The theory also runs for an algebra together with a derivation, and the
derivation then enters the complex. Ate\c{s}li, Esen and S\"utl\"u build such
a cohomology for an $n$-Lie algebra together with a representation
\cite[p.~2]{AtesliEsenSutlu2021}. Their complex carries the generalized
derivation extensions, and for $n=3$ it agrees with the known cohomology of
$3$-Lie algebras (Theorem~\ref{thm:cohomology:n-lie-derivation}) \cite[p.~1]{AtesliEsenSutlu2021}.

\begin{theorem}[{\cite[Proposition~2.13, pp.~25--26]{AtesliEsenSutlu2021}}]\label{thm:cohomology:n-lie-derivation}
Let $L$ be an $n$-Lie algebra and let $D\colon L^{n-1}\to L$ satisfy the
generalized derivation identity (2.10) of the cited paper, with
$D^{\sharp}\colon L^{n-2}\to\mathfrak{gl}(L)$ the map assigned to it. Then
$D^{\sharp}$ is a $1$-cocycle of $L$ with coefficients in
$\mathfrak{gl}(L)$.
\end{theorem}

So a generalized derivation is a cocycle in degree one, which is the
statement the paper aims at \cite[p.~2]{AtesliEsenSutlu2021}. The abelian
extensions and the infinitesimal deformations are read from the usual
cohomology of $n$-Lie algebras \cite[p.~1]{AtesliEsenSutlu2021}. The paper
also carries the Hochschild--Serre spectral sequence to $n$-Lie algebras
\cite[p.~1]{AtesliEsenSutlu2021}.

For $3$-Lie algebras Xu and Liu build the same theory from a $3$-LieDer pair,
a $3$-Lie algebra with a distinguished derivation
\cite[p.~1]{XuLiu2021ThreeLieDer}. Their complex controls the deformations
and the abelian extensions of the pair
\cite[p.~1]{XuLiu2021ThreeLieDer}. The degree $-1$ bracket on the complex
makes it a graded Lie algebra, and its Maurer--Cartan elements are exactly
the $3$-LieDer pairs \cite[p.~2]{XuLiu2021ThreeLieDer}.

A non-abelian cohomology also exists for such pairs. Das works with a
Nijenhuis Lie algebra, a Lie algebra with an operator $N$ that satisfies
$[N(x),N(y)]=N([N(x),y]+[x,N(y)]-N[x,y])$
\cite[Definition~2.1, p.~3]{Das2025Nijenhuis}. The non-abelian second group
of one such algebra with values in another classifies the non-abelian
extensions (Theorem~\ref{thm:cohomology:nijenhuis-nonabelian}) \cite[p.~2]{Das2025Nijenhuis}.

\begin{theorem}[{\cite[Theorem~3.4, p.~6]{Das2025Nijenhuis}}]\label{thm:cohomology:nijenhuis-nonabelian}
Let $(\mathfrak{g},N)$ and $(\mathfrak{h},S)$ be two Nijenhuis Lie algebras.
Then there is a bijection between the equivalence classes of non-abelian
extensions of $(\mathfrak{g},N)$ by $(\mathfrak{h},S)$ and the non-abelian
cohomology group $H^{2}_{\mathrm{nab}}((\mathfrak{g},N);(\mathfrak{h},S))$.
\end{theorem}

With abelian values the statement specializes to the usual second group
\cite[Theorem~3.7, p.~7]{Das2025Nijenhuis}. Das then turns to derivations.
Given an abelian extension of a Nijenhuis Lie algebra, a pair of Nijenhuis
Lie algebra derivations can be lifted exactly when it is compatible and its
Wells map vanishes \cite[Theorem~5.4, p.~19]{Das2025Nijenhuis}. The Wells map
here fits a short exact sequence \cite[Theorem~5.5, p.~20]{Das2025Nijenhuis}.

% Moved from chapters/20-applications/08-algebraic-topology.tex (source read in full: Weibel 1994, chapter 9).
